\documentclass[11pt]{article}
\usepackage{amsmath,amssymb,amsthm,fullpage}
\newtheorem{theorem}{Theorem}
\title{Independent Bits in Butterfly Rectangle Counts}
\author{}
\date{November 23, 2026}
\begin{document}
\maketitle

\paragraph{Status.}
This is a complete elementary written proof about permutation matrices,
checked exhaustively for the two- and three-layer examples. It is not a
HeadComplexity compiler: the required rectangle readouts have not been
implemented by admissible affine fractions with a saving. In particular,
this note does not prove either accepted asymptotic target. No original
model definition is changed. The theorem is not Lean checked.

\section{The independent rectangle system}
Let $d\geq1$, put $H=2^d$, and identify the wires $\{0,\ldots,H-1\}$
with bit strings $(x_0,\ldots,x_{d-1})$, where $x_0$ is the least
significant bit. For each $k=0,\ldots,d-1$, the $k$th layer contains
one binary switch for each pair
\[
 s\in\{0,1\}^k,\qquad g\in\{0,1\}^{d-k-1}.
\]
The switch indexed by $e=(k,s,g)$ either fixes or exchanges the two
wires $(s,0,g)$ and $(s,1,g)$. Different switches in a layer act on
disjoint pairs. Apply layers in increasing order of $k$. Given their
$dH/2$ binary settings $\eta_e$, their composition is a permutation
$\pi_\eta$ of the $H$ wires.

For each switch define sets of input and output labels
\[
 A_e=\{x:x_k=0,\ (x_{k+1},\ldots,x_{d-1})=g\},
 \qquad
 B_e=\{y:y_k=1,\ (y_0,\ldots,y_{k-1})=s\}.
\]
Both sets depend only on the switch index, not on any switch setting.
Their sizes are $2^k$ and $2^{d-k-1}$, respectively.

\begin{theorem}\label{thm:decoder}
For every switch setting $\eta$ and every switch $e$,
\[
 \bigl|\pi_\eta(A_e)\cap B_e\bigr|=\eta_e.
 \tag{1}\label{eq:bit}
\]
Consequently the $dH/2$ fixed rank-one rectangle functionals
\[
 L_e(P)=\sum_{x\in A_e}\sum_{y\in B_e}P_{xy}-\tfrac12
\]
on permutation matrices jointly realize all $2^{dH/2}$ strict sign
tables. Moreover, all $2^{dH/2}$ butterfly switch settings yield
different permutations.
\end{theorem}
\begin{proof}
Fix $e=(k,s,g)$. Every earlier layer changes only a bit with index
less than $k$. Hence it preserves the input set $A_e$, and its action
on that set is a bijection. Just before layer $k$, exactly one of the
inputs from $A_e$ occupies the wire $(s,0,g)$, and none occupies
$(s,1,g)$. More generally, for each lower tuple $s'$ exactly one
input from $A_e$ occupies $(s',0,g)$.

Layer $k$ keeps lower bits unchanged. Among the inputs from $A_e$,
the unique one with lower tuple $s$ therefore has its bit $k$ set
to one precisely when switch $e$ is exchanged, that is, precisely
when $\eta_e=1$. All later layers change only bits with index greater
than $k$, so preserve membership in $B_e$. This proves
\eqref{eq:bit}.

The rectangle coefficient matrix is the outer product
$\mathbf1_{A_e}\mathbf1_{B_e}^{\mathsf T}$ and thus has rank one.
Equation \eqref{eq:bit} gives $L_e(P_{\pi_\eta})=\eta_e-1/2$.
Every binary switch table is allowed independently, proving the
strict sign assertion. Equation \eqref{eq:bit} also recovers every
switch setting from the resulting permutation, proving injectivity.
\end{proof}

\paragraph{Exact verification.}
The companion script \texttt{butterfly-rectangle-decoder-check.py}
checks every switch setting for $d=1,2,3$, namely $2$, $16$, and
$4096$ settings. It checks every rectangle equality and permutation
injectivity. It also checks $1000$ deterministic pseudorandom
settings for $d=4$; that last test is not exhaustive. Its JSON
output is stored beside the script. These checks are independent
integer evaluations of the statement, not a substitute for the proof.

\section{The missing analytic implementation}
The theorem is stronger than counting permutations: it specifies
independent binary coordinates and explicit linear readouts. It also
differs from reading pairwise comparisons of permutation values.
Nevertheless, a single reciprocal matching perturbation typically
gives a readout of the form
\[
 \sum_x c_x V_e(\pi(x))K_e(\xi_x)
       -\text{a label-independent anchor term},
 \qquad K_e(\xi)=(t_e-\xi)^{-2},
\]
with all $t_e-\xi_x>0$ (and additional full-cube admissibility
conditions). A freely chosen color profile $V_e$ could supply
$\mathbf1_{B_e}$, but the remaining factor
$\mathbf1_{A_e}(x)$ is not a freely chosen row-dependent residue.
It would have to be provided by the kernel, a proved cancellation,
or another construction. Transposing the permutation exchanges
these two indicator roles.

The sets $A_e$ are dyadic intervals in ordinary binary order. The
sets $B_e$ are dyadic intervals after reversing bit order. Neither
family consists only of nested final intervals in a single ordering.
Exterior reciprocal kernels are not arbitrary interval indicators.
No simultaneous implementation of these indicators with sublinear
overhead has been proved here. Interpolating each indicator
separately with $H$ poles loses the proposed saving, and allowing
interior poles without checking denominator positivity changes the
model. The independent rectangle theorem supplies a concrete decoder
target; it does not close that analytic gap.

\section{A bound for literal exterior-kernel switch readouts}
The following restriction concerns a specific attempted implementation,
not the full admissible-head model. Fix positive integer $p$, token
positions $0\leq\xi_x\leq a$, and arbitrary signed residues $c_x$.
The tokens and residues are the same for every butterfly permutation.
A query may choose $t\geq b>a$, any profile $V$ on the output wires,
and any bias $\beta$, and reads
\[
 Q(\eta)=\beta+\sum_{x=0}^{H-1}
       \frac{c_x V(\pi_\eta(x))}{(t-\xi_x)^p}.
 \tag{2}\label{eq:literal}
\]
Say that this query reads a switch if its strict sign is that switch's
binary value, with positive sign for one, for every setting of all the
other switches. Different switches may use different queries.

\begin{theorem}\label{thm:firstlayer}
The number $R$ of first-layer switches readable by queries
\eqref{eq:literal} satisfies
\[
 R<(p+1)\left(1+\left(\frac b{b-a}\right)^{4p}\right).
 \tag{3}\label{eq:bound}
\]
The same bound holds if each queried switch may independently reverse
the sign convention. In particular it is independent of $H$ when
$a/b$ is bounded away from one.
\end{theorem}
\begin{proof}
Put $G=H/2$, and identify first-layer switches with pairs
$(2g,2g+1)$, $0\leq g<G$. Set
\[
 w_x(t)=\frac{c_x}{(t-\xi_x)^p},\qquad
 \Delta_g(t)=w_{2g}(t)-w_{2g+1}(t).
\]
Fix the first-layer settings and average the score over independent
uniform settings of every later switch. For any fixed input, its
output bit zero is unchanged by those layers. Each successive upper
bit becomes uniform when its own layer is applied, conditionally on
the preceding trajectory, and is never changed later. Consequently
the upper output bits are uniform conditional on output bit zero.
Let $v_0,v_1$ be the averages of $V$ on the two output parity classes.
The averaged score is therefore a constant plus
\[
 (v_1-v_0)\sum_{g=0}^{G-1}\Delta_g(t)\eta_{0,g}.
\]
If the original score strictly reads switch $g$, so does this finite
average, for every setting of the other first-layer switches. A
linear expression in independent bits can strictly read one bit
only if the magnitude of that bit's coefficient exceeds the sum of
the magnitudes of all its other coefficients: compare its minimum
when the target bit is one with its maximum when it is zero. Hence
\[
 |\Delta_g(t)|>\sum_{j\ne g}|\Delta_j(t)|.
 \tag{4}\label{eq:dominance}
\]
In particular $v_1-v_0$ and $\Delta_g(t)$ cannot vanish.

Select $R$ readable switches and their witness values $t_i$. For each
selected switch $j$ write
\[
 D_j(t)=(t-\xi_{2j})^p(t-\xi_{2j+1})^p,
 \qquad P_j(t)=D_j(t)\Delta_j(t).
\]
Every $P_j$ is a polynomial of degree at most $p$. The $R$ by $R$
matrix
\[
 L_{ij}=\frac{P_j(t_i)}{P_i(t_i)}
\]
therefore has rank at most $p+1$ and diagonal entries one. Moreover
\[
 \frac{D_j(t_i)}{D_i(t_i)}
 \leq\left(\frac{t_i}{t_i-a}\right)^{2p}
 \leq M:=\left(\frac b{b-a}\right)^{2p}.
\]
Using \eqref{eq:dominance}, even after omitting nonselected switches,
gives $\sum_{j\ne i}|L_{ij}|<M$. It follows that the squared
Frobenius norm satisfies $\|L\|_F^2<R(1+M^2)$.
For completeness, if $P$ is the orthogonal projection onto the
column space of $L$, then $PL=L$, and Frobenius Cauchy--Schwarz gives
\[
 |\operatorname{tr}L|=|\operatorname{tr}(PL)|
 \leq\|P\|_F\|L\|_F
 =\sqrt{\operatorname{rank}L}\,\|L\|_F.
\]
Since $\operatorname{tr}L=R$, these inequalities give
$R<(p+1)(1+M^2)$, as asserted.
\end{proof}

\paragraph{Scope.}
For the compact grouped transfer, $p=2$, $a<1/s$ and
$t_i=1-p_i$, with $\sum_i p_i<1$. At most one query has $t_i<1/2$.
Applying the theorem to all other queries yields
\[
 R<1+3\left(1+(1-2/s)^{-8}\right)\qquad(s>2).
\]
For instance, $s\geq100$ permits at most seven first-layer switch
queries. Thus the most literal realization of the full butterfly
rectangle decoder by one fixed-residue exterior kernel per query
cannot work at arbitrary size. This conclusion does not apply if
residues are allowed to depend on the complete switch setting, if
only a selected subfamily of butterfly permutations is used, if a
different code is used, or if additional admissible heads implement
new filters. It is not a lower bound for HeadComplexity.

\end{document}
